% ============================================
% LUC-03
% Gerado de molambudos.md — build_miolo.py
% ============================================
\fragtitle{luc}{L}{LUC}{03}{As Ruínas}
\markboth{LUC-03 --- As Ruínas}{}

\noindent \textit{O Colônia abandona}

\noindent \textbf{Março de 2026 --- Barbacena, Minas Gerais}

\vspace{0.5em}\noindent\rule{\textwidth}{0.5pt}\vspace{0.5em}

\noindent O Hospital Colônia de Barbacena foi fechado em 1980, um ano depois da morte de Joaquim.

\noindent Hoje, o que resta é um conjunto de prédios abandonados, cercados por mato alto e silêncio. As paredes ainda estão de pé, mas o telhado caiu em várias alas. O piso de madeira apodreceu. As camas de ferro enferrujado ainda estão lá, algumas com restos de colchão, outras com nada.

\noindent Lúcia chegou num sábado de manhã. O sol estava forte, mas o mato era alto demais para o sol alcançar o chão. Ela abriu caminho com as mãos, afastando galhos secos e teias de aranha.

\noindent A ala dos silenciados ficava nos fundos.

\noindent Ela reconheceu pela descrição no prontuário: \textit{\textquotedblleft{}Fim do corredor. Porta de ferro. Doze celas.\textquotedblright{}}

\noindent A porta de ferro estava entreaberta. Lúcia empurrou. O rangido ecoou pelo corredor vazio.

\noindent Dentro, o cheiro era de terra molhada e mofo. Por baixo, quase imperceptível, havia uma doçura de fruta estourada no fundo de uma lata.

\noindent Ela acendeu a lanterna do celular.

\noindent As celas eram pequenas. Dois metros por três. Em cada uma, uma cama de concreto. Em algumas, restos de cobertores --- trapos cinzentos que se desfizeram com o tempo.

\noindent A cela 1.260 ficava no final do corredor.

\noindent Lúcia parou na frente dela.

\noindent A porta era de madeira maciça, com uma abertura do tamanho de um prato --- o buraco por onde passavam a comida. Ela espiou.

\noindent Dentro, a cela estava vazia.
Quase.

\noindent No chão, no canto direito, havia um caderno.

\noindent O caderno não era o diário. O diário estava na mochila de Lúcia. Era outro caderno. Pequeno, de capa preta, com as páginas amareladas.

\noindent Ela abriu a porta. Entrou.

\noindent O cheiro adocicado era mais forte aqui.

\noindent Ela se abaixou e pegou o caderno. Na capa, escrito à mão:
\noindent A doçura podre se adensou ali, grudando na parte de trás da língua.

\noindent \textit{\textquotedblleft{}Joaquim --- 1945 a 1956\textquotedblright{}}

\noindent Os anos de silêncio.

\noindent Ela abriu. A caligrafia era minúscula, apertada, como se ele quisesse caber o máximo de palavras possível em cada página.

\noindent A primeira linha dizia:

\noindent \textquotedblleft{\textit{Ela não me deixa falar. Quando eu abro a boca, ela aperta minha garganta. A fome não quer que eu conte. Mas ela não pode me impedir de escrever. Escrever não é falar. Escrever é pensar no papel. Ela não lê meus pensamentos. Ela só lê o que eu escrevo. Mas ela está aprendendo.\textquotedblright{}}}

\noindent A última linha dizia:

\noindent \textit{\textquotedblleft{}1956. Hoje ela deixou eu falar. Perguntei por quê. Ela disse: 'Porque você já escreveu tudo. Agora a palavra está no papel. Não precisa mais sair pela boca.' Eu vou falar de novo. Mas antes vou escrever mais. Escrever é a única coisa que ainda é minha.\textquotedblright{}}

\noindent Lúcia fechou o caderno.

\noindent O cheiro adocicado estava mais forte.

\noindent Ela levantou os olhos.

\noindent No canto da cela, na parede, alguém tinha escrito alguma coisa. Com tinta. Ou com outra coisa.

\noindent \textquotedblleft{\textit{Bem-vinda, Dra. Lúcia. A senhora demorou.\textquotedblright{}}}

\noindent A caligrafia era a mesma do caderno.

\noindent A de Joaquim.

\noindent Escrita em 1956.

\noindent Prevendo que ela estaria ali. 70 anos depois.

\noindent Lúcia respirou fundo.

\noindent O cheiro adocicado encheu os pulmões.

\vspace{0.5em}
\noindent\textit{\textbf{↪ Rotas:} \rota{LUC-04} \rotasep \rota{MEM-08}}